\documentclass[11pt,a4paper]{report}
\usepackage{fontspec,amsmath,amssymb,amsthm,microtype,xcolor}
\setmainfont{lmroman10-regular.otf}[BoldFont=lmroman10-bold.otf,ItalicFont=lmroman10-italic.otf,BoldItalicFont=lmroman10-bolditalic.otf]
\usepackage[margin=25mm]{geometry}
\usepackage[unicode,hidelinks]{hyperref}
\hypersetup{pdfauthor={LangevinLean Project},bookmarksdepth=2}
\setlength{\emergencystretch}{3em}
\setlength{\parskip}{0.4em}
\setlength{\parindent}{0pt}
\setcounter{tocdepth}{1}
\makeatletter
\renewcommand*\l@section{\@dottedtocline{1}{1.5em}{3.8em}}
\renewcommand*\@pnumwidth{2.5em}
\renewcommand*\@tocrmarg{3.5em}
\makeatother
\setcounter{secnumdepth}{2}
\allowdisplaybreaks[2]
\raggedbottom
\title{\texorpdfstring{Interacting Particle Langevin under LSI}{Interacting Particle Langevin under LSI}\\[0.6em]\large Mathematical proofs}
\author{LangevinLean Project}
\date{Draft report LL-04\\[0.4em]\small Content revision 4eecf839e463}
\begin{document}
\pagenumbering{roman}
\maketitle
\chapter*{About this document}
\addcontentsline{toc}{chapter}{About this document}
These draft notes present the assumptions, convergence theorems and detailed mathematical proofs for this development. The supporting arguments are arranged so that prerequisites precede their applications. They use ordinary mathematical notation throughout.

Prepared by the LangevinLean Project with AI assistance, these notes form part of an ongoing effort to formalize Langevin algorithms comprehensively. This is draft work and may contain mistakes.

\section*{Literature and scope}
The particle entropy argument presented here is unpublished; the cited works provide its sampling and functional-inequality background.

\begin{enumerate}
\item Santosh S. Vempala and Andre Wibisono (2019). \emph{Rapid Convergence of the Unadjusted Langevin Algorithm: Isoperimetry Suffices}. Advances in Neural Information Processing Systems 32, 8094–8106. \url{https://proceedings.neurips.cc/paper/2019/hash/65a99bb7a3115fdede20da98b08a370f-Abstract.html}.
\item Alex Glyn-Davies, Connor Duffin, Ieva Kazlauskaite, Mark Girolami, and Ö. Deniz Akyildiz (2025). \emph{Statistical Finite Elements via Interacting Particle Langevin Dynamics}. SIAM/\allowbreak{}ASA Journal on Uncertainty Quantification 13(3), 1200–1227. \url{https://doi.org/10.1137/24M1693593}.
\end{enumerate}
These proofs use results established in the earlier project reports \cite{LangevinLean2026ULASC,LangevinLean2026ULALSI,LangevinLean2026IPLA}. Each application below gives the section containing the required result. The earlier proofs are not reproduced in this report.

\tableofcontents
\clearpage
\pagenumbering{arabic}
\chapter{Assumptions and main results}
\section{\texorpdfstring{Joint KL bound and the linear lifted-\allowbreak{}LSI specialization}{Joint KL bound and the linear lifted-LSI specialization}}


Let \(N\ge1\), \(D=d_\theta+Nd_x\), assume \(D\ge1\), and let the joint potential \(U(\theta,x)\) have an \(L\)-Lipschitz gradient with \(L>0\). In scaled coordinates set


\[z=(\theta,x^1/\sqrt N,\ldots,x^N/\sqrt N),\qquad
V_N(z)=\sum_{i=1}^NU(\theta,\sqrt N z_i),\quad
\Pi_N(dz)\propto e^{-V_N(z)}dz.\]


The target is the Gibbs law of the lifted potential under the established target identity \cite[Section~3.5.1]{LangevinLean2026IPLA}. Assume an LSI with constant \(\beta_N>0\) for this actual lifted target. The positive-dimensional compatibility result \cite[Section~3.6.9]{LangevinLean2026ULALSI} gives \(\beta_N\le NL\); the final theorem derives this compatibility from its stated assumptions. For finite initial joint entropy and


\[0<\gamma\le\frac{\beta_N}{4NL^2},\]


the IPLA laws \(\mu_k\), equivalently ULA for \(V_N\) at step \(\gamma/N\) by the established algorithm-law identity \cite[Section~3.16.3]{LangevinLean2026IPLA}, satisfy


\[\frac{H_{\Pi_N}(\mu_k)}N\le
e^{-\beta_N\gamma k/N}\frac{H_{\Pi_N}(\mu_0)}N
+\frac{8\gamma DL^2}{\beta_N}.\]


In particular, if the lifted target has the admissible constant \(\beta_N=N\alpha\), then


\[\boxed{\frac{H_{\Pi_N}(\mu_k)}N\le
e^{-\alpha\gamma k}\frac{H_{\Pi_N}(\mu_0)}N
+\frac{8\gamma L^2}{\alpha}\left(d_x+\frac{d_\theta}N\right).}\]


The condition is now \(0<\gamma\le\alpha/(4L^2)\). An LSI for the original pair target is not substituted for the lifted-target hypothesis. The result derives all interpolation estimates from the shared ULA entropy theorem \cite[Section~3.29.7]{LangevinLean2026ULALSI} and finite initial joint KL. If \(D=0\), both laws are the unique point mass and their entropy is zero.


{\small\color{gray}Supporting proof: Section~\ref{proof-module-IPLA-LSI-UnconditionalTheorem-lean}.\par}
\chapter{Mathematical argument}
\section{\texorpdfstring{The geometry fixes the dimension, gradient constant, and step}{The geometry fixes the dimension, gradient constant, and step}}


Each physical pair \((\theta,x^i)\) is extracted from the scaled state by \((\theta,\sqrt N z_i)\). The two gradient-block formulas \cite[Section~3.1.5]{LangevinLean2026IPLA} show that differentiating the lifted sum gives the averaged parameter drift and ordinary latent drifts at step \(\gamma/N\). The corresponding Euclidean state has


\[D=d_\theta+Nd_x,\qquad L_N=NL,\qquad h_N=\gamma/N.\]


The dimension identity is the coordinate-space dimension formula \cite[Section~3.1.1]{LangevinLean2026IPLA}. The parameter block of the lifted gradient is a sum over particles; each latent block contributes a factor \(\sqrt N\). Squaring and summing the blocks, then applying Cauchy–Schwarz to the common-parameter sum, proves the gradient Lipschitz constant \(NL\) \cite[Section~3.1.4]{LangevinLean2026IPLA}. The coordinate Gaussian laws and independence identify the actual noise-driven IPLA law with this ULA transition \cite[Section~3.16.3]{LangevinLean2026IPLA}. No equality between \(d_\theta\) and \(d_x\) is needed.



The scaling can also be verified directly in physical coordinates. The parameter component of a step is the previously established physical parameter update \cite[Section~3.1.7]{LangevinLean2026IPLA}:

\[\theta_{k+1}=\theta_k-\frac\gamma N\sum_i\nabla_\theta U(\theta_k,x_k^i)
+\sqrt{\frac{2\gamma}{N}}Z_k^\theta.\]

For the latent block, multiply the scaled update by \(\sqrt N\); the physical latent update \cite[Section~3.1.7]{LangevinLean2026IPLA} is

\[x_{k+1}^i=x_k^i-\gamma\nabla_xU(\theta_k,x_k^i)+\sqrt{2\gamma}Z_k^i.\]

Thus the lifted Gaussian transition is exactly the intended particle algorithm, with the same joint innovation law. The dimension \(D\) counts all coordinates of that joint state, and the full joint entropy is the natural input to the shared ULA entropy theorem \cite[Section~3.29.7]{LangevinLean2026ULALSI}.

\section{\texorpdfstring{Apply the ULA one-\allowbreak{}step estimate with the lifted target}{Apply the ULA one-step estimate with the lifted target}}


Apply the ULA one-step entropy theorem \cite[Section~3.29.6]{LangevinLean2026ULALSI} to \((V_N,\Pi_N)\), with constant \(\beta_N\), dimension \(D\), smoothness \(NL\), and step \(\gamma/N\). Its step restriction becomes


\[\frac\gamma N\le\frac{\beta_N}{4(NL)^2}
\quad\Longleftrightarrow\quad
\gamma\le\frac{\beta_N}{4NL^2}.\]


The one-step forcing simplifies without losing a particle factor:


\[6h_N^2D L_N^2=6(\gamma/N)^2D(NL)^2=6\gamma^2DL^2.\]


Consequently the unnormalized joint entropy obeys


\[H_{k+1}\le e^{-\beta_N\gamma/N}H_k+6\gamma^2DL^2.\]


The lifted LSI assumption includes probability normalization of the Gibbs target \cite[Section~3.1.3]{LangevinLean2026ULALSI}. The shared finiteness theorem \cite[Section~3.29.2]{LangevinLean2026ULALSI} supplies finite KL before the one-step theorem is applied and its rescaled constants are simplified. Every analytic input is supplied by the actual Gaussian ULA construction.



The product \(\beta_Nh_N\) equals \(\beta_N\gamma/N\), which fixes the rate in the one-step exponential. The numerical step condition also yields \(\beta_N\gamma/N\le1/4\) through the lifted compatibility inequality \cite[Section~3.6.9]{LangevinLean2026ULALSI}. Hence the geometric-denominator estimate \cite[Section~3.20.2]{LangevinLean2026ULALSI} used in the ULA proof is available with the rescaled constants, without losing a particle factor.



The recurrence is for the unnormalized joint divergence. The additive term has no remaining factor \(N^{-1}\) because the squared step contributes \(N^{-2}\) and the squared lifted smoothness contributes \(N^2\). Normalization by \(N\) is postponed until the complete joint inequality has been derived.

\section{\texorpdfstring{Normalize by the particle number only after deriving the joint bound}{Normalize by the particle number only after deriving the joint bound}}


The shared full ULA entropy theorem \cite[Section~3.29.7]{LangevinLean2026ULALSI} gives the unnormalized floor after the lifted substitution:


\[\frac{8h_NDL_N^2}{\beta_N}
=\frac{8\gamma NDL^2}{\beta_N}.\]


Divide the entire inequality by positive \(N\), including the initial entropy and the floor. This proves


\[\frac{H_k}N\le e^{-\beta_N\gamma k/N}\frac{H_0}N
+\frac{8\gamma DL^2}{\beta_N}.\]


Now, and only under the separate lifted-target assumption \(\beta_N=N\alpha\), use \(D/N=d_x+d_\theta/N\). The rate is \(e^{-\alpha\gamma k}\), and the per-particle floor is \((8\gamma L^2/\alpha)(d_x+d_\theta/N)\). Shared parameter coordinates prevent a direct tensorization argument from a pair LSI; the target for which LSI is assumed remains \(\Pi_N\).



The normalization is a scalar division of an already established inequality, not a change of reference measure. Under \(\beta_N=N\alpha\), substitute \(D=d_\theta+Nd_x\) only after the division:

\[\frac{8\gamma DL^2}{N\alpha}
=\frac{8\gamma L^2}{\alpha}\left(\frac{d_\theta}{N}+d_x\right).\]

This exposes two distinct dimension contributions. The shared parameter dimension is amortized over particles, while the latent dimension contributes once per particle and survives in the normalized floor. The assumption \(\beta_N=N\alpha\) concerns the interacting lifted target; shared coordinates mean it cannot be inferred by independent tensorization of unrelated pair laws.

\section{\texorpdfstring{Project to the parameter marginal using data processing}{Project to the parameter marginal using data processing}}


Let \(P_\theta\) be the continuous linear parameter projection \cite[Section~3.1.6]{LangevinLean2026IPLA}, \(\mu_k^\theta=P_{\theta\#}\mu_k\), and \(\Pi_N^\theta=P_{\theta\#}\Pi_N\). Relative entropy decreases under a measurable map:


\[H_{\Pi_N^\theta}(\mu_k^\theta)\le H_{\Pi_N}(\mu_k).\]


This controls the unnormalized parameter KL. In particular, under \(\beta_N=N\alpha\),


\[H_{\Pi_N^\theta}(\mu_k^\theta)
\le e^{-\alpha\gamma k}H_{\Pi_N}(\mu_0)
+\frac{8\gamma NL^2}{\alpha}\left(d_x+\frac{d_\theta}N\right).\]


The initial entropy remains the full joint entropy. Data processing does not allow it to be replaced by the smaller initial parameter entropy on the right. Likewise, a normalized joint bound must be multiplied by \(N\) before it bounds the unscaled parameter KL. The proof establishes finite joint KL, transfers finiteness to the marginal, and then compares their ordinary finite values.



For a density-based view of data processing, let \(R=d\mu_k/d\Pi_N\). The density of the parameter marginal is the conditional mean \(\mathbb E_{\Pi_N}[R\mid\theta]\). Convexity of \(u\log u\) and conditional Jensen therefore give

\[\int\mathbb E[R\mid\theta]\log\mathbb E[R\mid\theta]d\Pi_N
\le\int R\log R\,d\Pi_N.\]

This proves the marginal inequality with the original, unnormalized joint entropy on the right. There is no factor \(N^{-1}\) in data processing itself. A bound for entropy per particle must consequently be multiplied by \(N\) before being used as a bound on the ordinary parameter marginal KL.

\chapter{Supporting proofs}
We now supply the detailed supporting arguments. Each section introduces the objects and estimates needed by the results that follow it.

\section{\texorpdfstring{Exact entropy scaling with an explicit lifted LSI constant}{Exact entropy scaling with an explicit lifted LSI constant}}
\label{proof-module-IPLA-LSI-Scaling-lean}


Relative entropy is measured first on the full enlarged state. Dividing it by particle count produces a separate intensive quantity. The calculations below track both quantities, including their initial terms and forcing constants, under the physical-time substitution. The actual enlarged-target logarithmic Sobolev constant is retained throughout rather than inferred from an unrelated one-pair target.




Keep the LSI constant \(\beta\) of the actual lifted target explicit. Substitute \(h=\gamma/N\), \(\ell=NL\), \(D=d_\theta+Nd_x\) into the shared ULA entropy argument. The regime \(\beta=N\alpha\) is an additional lifted-target assumption, with its own scalar specializations.

\subsection{\texorpdfstring{Translate the numerical step and compatibility conditions}{Translate the numerical step and compatibility conditions}}
\label{report-step-IPLA-LSI-Scaling-lean-1}


For \(N>0\), \(L>0\), factor the denominator and cancel positive \(N\) to get

\[\frac\gamma N\le\frac\beta{4(NL)^2}\quad\Longleftrightarrow\quad\gamma\le\frac\beta{4NL^2}.\]

If \(\beta=N\alpha\), multiplication by positive \(N\) gives \(N\alpha\le NL\Longleftrightarrow\alpha\le L\), and cancellation gives \(N\alpha/(4NL^2)=\alpha/(4L^2)\). The forward compatibility lemma simply uses the equivalence. These are exact ordered-field identities; they do not derive an LSI for the lifted model from any pair-model LSI.




The general lifted theorem uses \(h=\gamma/N\), score constant \(L_N=NL\), dimension \(D=d_\theta+Nd_x\), and an LSI constant \(\beta\) for the actual lifted Gibbs target. Thus \(h\le\beta/(4L_N^2)\) becomes \(\gamma\le\beta/(4NL^2)\). Only after additionally assuming \(\beta=N\alpha\) can the step restriction simplify to \(\gamma\le\alpha/(4L^2)\). Positive \(N\) justifies these cancellations; no product tensorization is used, since the pair potentials share a parameter.





The smoothness constant is squared in the entropy step restriction, so lifting introduces \(N^2\) in its denominator. The Euler step contributes one compensating factor \(1/N\); the physical restriction consequently retains one factor \(N\). If the target's LSI constant also grows as \(N\), that last factor cancels. These are numerical equivalences under positive denominators and should be separated from proving any functional inequality.

\subsection{\texorpdfstring{Rewrite the contraction exponent}{Rewrite the contraction exponent}}
\label{report-step-IPLA-LSI-Scaling-lean-2}


Associativity of real multiplication and division gives \(-\beta(\gamma/N)k=-\beta\gamma k/N\), so the corresponding exponentials agree. When \(\beta=N\alpha\) and \(N>0\), cancellation yields

\[e^{-\beta\gamma k/N}=e^{-\alpha\gamma k},\qquad e^{-(N\alpha)(\gamma/N)k}=e^{-\alpha\gamma k}.\]

The second version composes the first two identities. The rate becomes independent of \(N\) only under this explicitly supplied linear scaling of the actual lifted LSI constant.





At one iteration, the damping exponent is \(\beta h\). After \(k\) iterations it is \(\beta hk\), corresponding to the elapsed lifted time. Physical time is \(k\gamma\), so the effective entropy decay per physical unit is \(\beta/N\). Only a supplied lower bound of order \(N\) for the enlarged-target LSI constant makes this rate uniform in particle count.

\subsection{\texorpdfstring{Rewrite full and normalized entropy biases}{Rewrite full and normalized entropy biases}}
\label{report-step-IPLA-LSI-Scaling-lean-3}


The shared ULA bias is \(8hD\ell^2/\beta\). Substitute the lifted values and cancel one particle factor:

\[\frac{8(\gamma/N)D(NL)^2}{\beta}=\frac{8\gamma ND L^2}{\beta},\qquad \frac1N\frac{8\gamma ND L^2}{\beta}=\frac{8\gamma DL^2}{\beta}.\]

The dimension identity \(D/N=d_x+d_\theta/N\) follows by distributing division over the sum. If \(\beta=N\alpha\), the normalized bias becomes

\[\frac{8\gamma DL^2}{N\alpha}=\frac{8\gamma L^2}{\alpha}\left(d_x+\frac{d_\theta}{N}\right).\]

The formal algebra handles \(\alpha=0\) separately because real division is total; applications use \(\alpha>0\). For any bound \(H\le qH_0+8\gamma ND L^2/\beta\), divide by positive \(N\), distribute, and use the second identity to obtain the normalized entropy bound.




The full-state bias is \(8hDL_N^2/\beta=8\gamma NDL^2/\beta\); dividing entropy by \(N\) removes exactly one particle factor. Under \(\beta=N\alpha\), its intensive form becomes \((8\gamma L^2/\alpha)(d_x+d_\theta/N)\). The initial entropy is divided by \(N\) at the same time. The one-step forcing is different: \(6h^2DL_N^2=6\gamma^2DL^2\), so its normalized form is \(6\gamma^2L^2(d_x+d_\theta/N)\). Keeping these full and intensive quantities separate avoids losing a factor of \(N\).





Normalization divides every term by the same positive \(N\). In particular it does not change the exponential factor and cannot leave the full initial entropy unchanged. The bias also contains \(D\), which itself grows with particle count. Rewriting \(D/N\) reveals the intensive dimension and permits a meaningful comparison of bounds at different particle counts under the additional linear-LSI assumption.

\subsection{\texorpdfstring{Rewrite the one-\allowbreak{}step forcing and normalized recurrence}{Rewrite the one-step forcing and normalized recurrence}}
\label{report-step-IPLA-LSI-Scaling-lean-4}


The squared step cancels both particle factors in the squared smoothness:

\[6(\gamma/N)^2D(NL)^2=6\gamma^2DL^2.\]

Dividing by \(N\) and substituting the intensive dimension gives \(6\gamma^2DL^2/N=6\gamma^2L^2(d_x+d_\theta/N)\). Thus any supplied lifted recurrence normalizes to

\[\frac{H_{k+1}}N\le e^{-\beta\gamma/N}\frac{H_k}N+\frac{6\gamma^2DL^2}{N}.\]

This scalar helper explicitly assumes its recurrence. Its proof divides by positive \(N\), rewrites the exponential argument, and applies the forcing identity. Under \(\beta=N\alpha\), the rate and intensive-dimension identities give the specialized recurrence with rate \(e^{-\alpha\gamma}\) and forcing \(6\gamma^2L^2(d_x+d_\theta/N)\). The analytic theorem derives the recurrence separately from interpolation data.





The forcing belongs to a one-step recurrence and is quadratic in the step. The long-time floor is larger by the geometric denominator \(1-e^{-\beta h}\), of order \(\beta h\). This explains why the final bias is first order in the physical step even though the one-step error is second order. Scalar normalization preserves the recurrence; the analytic interpolation argument is what supplies it for actual laws.

\section{\texorpdfstring{Conditional entropy convergence for the actual Gaussian particle laws}{Conditional entropy convergence for the actual Gaussian particle laws}}
\label{proof-module-IPLA-LSI-Theorem-lean}


For the conditional entropy route, the actual Euler interpolation is assumed to have the density, weak-equation, differentiability and integrability properties needed for entropy calculus. These explicit analytic assumptions connect density calculations with the actual Gaussian algorithm laws. The proof derives finite entropy, the one-step recurrence and its geometric summation from those properties and the enlarged-target LSI.




Earlier results used here: \cite[Section~3.16]{LangevinLean2026IPLA}; \cite[Section~3.5]{LangevinLean2026IPLA}; \cite[Section~3.34]{LangevinLean2026ULALSI}.
\subsection{\texorpdfstring{The canonical particle laws are the actual Gaussian algorithm laws}{The canonical particle laws are the actual Gaussian algorithm laws}}
\label{report-step-IPLA-LSI-Theorem-lean-1}




The Gaussian transition is a measurable image of the product of the current law and one independent standard vector. Its total mass is therefore one. Applying it repeatedly defines genuine probability laws without any stationarity assumption. Independence of the supplied innovation path identifies these same laws with the explicit particle recursion, so the entropy estimate concerns its current distribution at each finite iteration.



Earlier results used here: \cite[Section~3.8.4]{LangevinLean2026ULALSI}; \cite[Section~3.8.6]{LangevinLean2026ULALSI}; \cite[Section~3.1.1]{LangevinLean2026IPLA}; \cite[Section~3.1.4]{LangevinLean2026IPLA}; \cite[Section~3.16.3]{LangevinLean2026IPLA}; \cite[Section~3.3.2]{LangevinLean2026IPLA}.
\subsection{\texorpdfstring{Density regularity proves finite actual KL before real conversion}{Density regularity proves finite actual KL before real conversion}}
\label{report-step-IPLA-LSI-Theorem-lean-2}


The LSI assumption includes that \(\Pi_N\) is a probability measure. The interpolation data give positive relative densities \(r_k(t)\) with \(r_k(t)\Pi_N\) equal to the actual frozen-drift interpolation law, together with local entropy regularity and spatial smoothness. The shared finiteness theorem applies with \(h=\gamma/N\ge0\) and proves

\[\operatorname{KL}(\rho_k\Vert\Pi_N)<\infty\quad\text{for every }k.\]



For a relative density \(r\), entropy is \(\int r\log r\,d\Pi_N\) when this quantity is integrable. The assumptions ensure that the interpolation density represents the Euler law at the relevant interval endpoint and that its entropy is finite. Equality of endpoint measures then transfers this finiteness to \(\rho_k\). This step is logically prior to treating the divergence as an ordinary real number in a differential or scalar recurrence.



Earlier results used here: \cite[Section~3.26.3]{LangevinLean2026ULASC}; \cite[Section~3.1.3]{LangevinLean2026ULALSI}; \cite[Section~3.34.1]{LangevinLean2026ULALSI}; \cite[Section~3.1.1]{LangevinLean2026IPLA}; \cite[Section~3.3.2]{LangevinLean2026IPLA}.
\subsection{\texorpdfstring{Derive entropy production and the one-\allowbreak{}step inequality}{Derive entropy production and the one-step inequality}}
\label{report-step-IPLA-LSI-Theorem-lean-3}


Put \(h=\gamma/N\), \(\ell=NL\), and \(H_k=\operatorname{KL}(\rho_k\Vert\Pi_N)\). The data provide the weak interpolation equation, differentiation and integration regularity, an error field \(e_k(t)\), and the bounds

\[\int r_k(t)\|e_k(t)\|^2d\Pi_N\le t^2\ell^2G_k+2tD\ell^2,\qquad G_k\le\frac{4\ell^2}{\beta}H_k+2D\ell.\]

The shared calculus tests the weak equation against the logarithmic density and uses Young's inequality to derive \(H'(t)\le-3J(t)/4+\int r\|e\|^2d\Pi_N\). LSI gives \(J(t)\ge2\beta H(t)\). Since \(h\le\beta/(4\ell^2)\) and \(\beta\le\ell\), one has \(h\ell\le1/4\) and \(\beta h\le1/4\). The error bounds then imply

\[H'(t)\le-\frac32\beta H(t)+\frac\beta4H_k+\frac52hD\ell^2.\]

The shared affine differential comparison and \(e^{-3x/2}+x/4\le e^{-x}\) for \(0\le x\le1/4\) give \(H_{k+1}\le e^{-\beta h}H_k+6h^2D\ell^2\). The IPLA proof checks positivity of \(L\) from \(0<\beta\le NL\), translates the step restriction, and substitutes the exact lifted dimension and forcing. Therefore

\[H_{k+1}\le e^{-\beta\gamma/N}H_k+6\gamma^2DL^2.\]

Dividing by positive \(N\) gives the normalized recurrence with forcing \(6\gamma^2L^2(d_x+d_\theta/N)\). Neither recurrence is a field of the interpolation structure.






The error field measures the difference between the true Langevin drift and the frozen Euler drift over one interval. Young's inequality absorbs part of its pairing with the density score into Fisher information. The LSI converts the remaining Fisher-information dissipation into entropy decay. The step restriction ensures the drift-error contribution proportional to the old entropy is small enough to preserve a strict contraction, leaving the explicit dimension-dependent forcing.



Earlier results used here: \cite[Section~3.1.3]{LangevinLean2026ULALSI}; \cite[Section~3.34.1]{LangevinLean2026ULALSI}; \cite[Section~3.1.1]{LangevinLean2026IPLA}; \cite[Section~3.1.3]{LangevinLean2026IPLA}; \cite[Section~3.1.4]{LangevinLean2026IPLA}; \cite[Section~3.3.2]{LangevinLean2026IPLA}.
\subsection{\texorpdfstring{Iterate and simplify the finite-\allowbreak{}entropy bound}{Iterate and simplify the finite-entropy bound}}
\label{report-step-IPLA-LSI-Theorem-lean-4}


Let \(q=e^{-\beta h}\). Repeated substitution in the one-step estimate gives

\[H_k\le q^kH_0+6h^2D\ell^2\sum_{j=0}^{k-1}q^j.\]

Because \(0<\beta h\le1/4\), the elementary bound \(1-e^{-\beta h}\ge3\beta h/4\) bounds the geometric sum by \(4/(3\beta h)\). Thus the shared theorem proves \(H_k\le e^{-\beta hk}H_0+8hD\ell^2/\beta\). Substitute \(h=\gamma/N\), \(\ell=NL\), and \(D=d_\theta+Nd_x\) to obtain

\[H_k\le e^{-\beta\gamma k/N}H_0+\frac{8\gamma ND L^2}{\beta}.\]

The proof invokes the shared bound after deriving every lifted numerical and regularity input, then uses the exact scaling identities. Dividing the entire inequality by positive \(N\) yields

\[\frac{H_k}{N}\le e^{-\beta\gamma k/N}\frac{H_0}{N}+\frac{8\gamma DL^2}{\beta}.\]


Iterating the finite-entropy inequality uses the shared ULA entropy bound, whose geometric sum is controlled with the positive exponent \(\beta h\). After the exact lifted substitutions it yields \[H_k\le e^{-\beta\gamma k/N}H_0+\frac{8\gamma NDL^2}{\beta}.\] The normalized result divides both the initial term and bias by \(N\). The linear-LSI specialization then substitutes the additional lifted hypothesis \(\operatorname{LSI}(\Pi_N,N\alpha)\), producing rate \(e^{-\alpha\gamma k}\) and intensive dimension \(d_x+d_\theta/N\). The Gaussian realization corollaries use equality of laws; they introduce no new entropy inequality.





Repeated substitution gives a finite geometric series with a nonnegative ratio strictly smaller than one. Its denominator is positive because \(\beta,h>0\). The elementary lower bound on \(1-e^{-\beta h}\) converts the quadratic forcing to the displayed linear-step floor. No passage to an infinite-time limit is needed: the same floor bounds the finite sum for every finite \(k\), while the initial contribution retains its exact exponential factor.



Earlier results used here: \cite[Section~3.1.3]{LangevinLean2026ULALSI}; \cite[Section~3.34.1]{LangevinLean2026ULALSI}; \cite[Section~3.1.1]{LangevinLean2026IPLA}; \cite[Section~3.1.3]{LangevinLean2026IPLA}; \cite[Section~3.1.4]{LangevinLean2026IPLA}; \cite[Section~3.3.2]{LangevinLean2026IPLA}.
\subsection{\texorpdfstring{Specialize only when the lifted target supplies \ensuremath{\beta} = N\ensuremath{\alpha}}{Specialize only when the lifted target supplies β = Nα}}
\label{report-step-IPLA-LSI-Theorem-lean-5}


Assume the actual lifted target has the supplied LSI constant \(N\alpha\), with \(0<\alpha\le L\), and \(0<\gamma\le\alpha/(4L^2)\). Then \(N\alpha>0\), compatibility \(N\alpha\le NL\) holds, and the step equivalence supplies the general theorem's restriction. Its rate cancels to \(e^{-\alpha\gamma k}\); the normalized bias uses \(D/N=d_x+d_\theta/N\). Consequently

\[\frac{H_k}{N}\le e^{-\alpha\gamma k}\frac{H_0}{N}+\frac{8\gamma L^2}{\alpha}\left(d_x+\frac{d_\theta}{N}\right).\]

This is a specialization of a lifted-target LSI already given in the hypotheses. The shared parameter makes the lifted model nonproduct, so no pair-target LSI tensorization is inferred.





Setting \(\beta=N\alpha\) in the general theorem changes both the admissible physical step and the decay rate. Dividing by \(N\) then expresses the floor using \(d_x+d_\theta/N\). The shared parameter prevents treating the enlarged Gibbs distribution as independent pair targets, so linear growth of its LSI strength remains a stated premise. The resulting specialization is valid whenever that premise is actually available.



Earlier results used here: \cite[Section~3.1.3]{LangevinLean2026ULALSI}; \cite[Section~3.34.1]{LangevinLean2026ULALSI}; \cite[Section~3.1.1]{LangevinLean2026IPLA}; \cite[Section~3.3.2]{LangevinLean2026IPLA}.
\subsection{\texorpdfstring{Transfer the canonical bounds to the supplied Gaussian realization}{Transfer the canonical bounds to the supplied Gaussian realization}}
\label{report-step-IPLA-LSI-Theorem-lean-6}


For a square-integrable initialization on \(P\) and noises satisfying all scalar Gaussian and independence fields, the law identity gives \(\rho_k=\operatorname{laws}(P,U,\gamma,k)\). Rewriting the normalized general bound by this equality proves the first Gaussian realization corollary. Rewriting the specialized \(N\alpha\)-LSI bound proves the second one, with all algorithm and interpolation assumptions retained. The entropy on the right is the full initial joint divergence relative to \(\Pi_N\); normalization divides both it and the subsequent joint divergence by \(N\). These conclusions by themselves do not assert parameter concentration about a mode.





Relative entropy depends only on two measures, so equality of output laws is sufficient to transfer every conclusion. The initial joint law is also the same under this identification. Thus neither its entropy nor its normalization changes when replacing the canonical product realization by a supplied Gaussian realization. The interpolation assumptions remain attached to the same Euler laws, and no extra entropy inequality is introduced by this transfer.



Earlier results used here: \cite[Section~3.1.3]{LangevinLean2026ULALSI}; \cite[Section~3.34.1]{LangevinLean2026ULALSI}; \cite[Section~3.1.1]{LangevinLean2026IPLA}; \cite[Section~3.3.2]{LangevinLean2026IPLA}.
\section{\texorpdfstring{Parameter KL by data processing and guarded real conversion}{Parameter KL by data processing and guarded real conversion}}
\label{proof-module-IPLA-LSI-Marginal-lean}


For entropy, projecting onto the parameter block requires only measurable data processing. No explicit formula for the parameter density is needed. This route therefore applies under the enlarged-target LSI assumptions even when no strong-convexity argument or concentration about a likelihood maximizer is available.




The parameter reference law is exactly the projection of the actual lifted Gibbs target. Measurable projection decreases extended KL, and finite joint KL makes the real bound transfer valid. This measure-theoretic argument needs no marginal-density or strong-convexity formula.



Earlier results used here: \cite[Section~3.3]{LangevinLean2026IPLA}.
\subsection{\texorpdfstring{Define the projected target and prove its probability property}{Define the projected target and prove its probability property}}
\label{report-step-IPLA-LSI-Marginal-lean-1}


With \(p=p_\theta\), define \(\Pi_N^\theta=p_\#\Pi_N\). The projection is a continuous linear map, hence measurable. For any probability measure \(\mu\), \((p_\#\mu)(\mathbb R^{d_\theta})=\mu(p^{-1}(\mathbb R^{d_\theta}))=\mu(\mathbb R^D)=1\). This proves probability preservation for any state law and then for the named target whenever \(\Pi_N\) is a probability. In the conditional LSI application that normalization comes from the LSI assumption itself.





Pushforward preserves probability because the inverse image of the whole target space is the whole source space. The definition \(\Pi_N^\theta=(p_\theta)_\#\Pi_N\) therefore fixes a genuine parameter reference law whenever the joint target is normalized. It need not be described by a marginal potential for the entropy argument. The same observation applies to the current particle-state law and its parameter projection.



Earlier results used here: \cite[Section~3.1.1]{LangevinLean2026IPLA}; \cite[Section~3.1.6]{LangevinLean2026IPLA}; \cite[Section~3.3.2]{LangevinLean2026IPLA}.
\subsection{\texorpdfstring{Extended KL decreases under the parameter projection}{Extended KL decreases under the parameter projection}}
\label{report-step-IPLA-LSI-Marginal-lean-2}


The shared data-processing theorem for a measurable map yields

\[\operatorname{KL}(p_\#\mu\Vert p_\#\nu)\le\operatorname{KL}(\mu\Vert\nu).\]

At the density level this is convexity of \(u\log u\) under averaging over forgotten coordinates; the measure theorem does not require a constructed conditional density. If the right-hand divergence is finite, the inequality immediately excludes infinity on the left. Since both extended quantities are finite, order is preserved when taking their real values. This proves the real data-processing inequality with its necessary finiteness guard.






Data processing averages out information in the discarded particle coordinates. Convexity of the entropy integrand makes this averaging nonincreasing for divergence. The resulting extended inequality remains meaningful if the joint divergence is infinite. For a finite numerical estimate, the joint entropy theorem first excludes that case; the projected divergence is then finite as well and both sides can be read as real numbers.



Earlier results used here: \cite[Section~3.1.1]{LangevinLean2026IPLA}; \cite[Section~3.1.6]{LangevinLean2026IPLA}.
\subsection{\texorpdfstring{Transfer a full-\allowbreak{}joint or per-\allowbreak{}particle bound with the correct factor}{Transfer a full-joint or per-particle bound with the correct factor}}
\label{report-step-IPLA-LSI-Marginal-lean-3}


If finite full joint KL has real bound \(H\le B\), transitivity with real data processing gives parameter KL at most \(B\). If instead \(H/N\le B\) and \(N>0\), multiplication by the positive count yields \(H\le NB\), hence

\[\operatorname{KL}(p_\#\mu\Vert p_\#\nu)\le NB.\]

The proof uses the positive-denominator order equivalence before applying the full-joint transfer lemma. This factor is necessary because normalized joint entropy is not the same quantity as parameter-marginal entropy.




If the available estimate is intensive, \(H/N\le B\), positive \(N\) gives the full joint bound \(H\le NB\), and data processing then gives marginal KL at most \(NB\). It does not give marginal KL at most \(B\). The initial entropy in the final bound is still full joint entropy: projection decreases initial entropy, so replacing that term by initial parameter entropy would make the right-hand side smaller without a supporting argument.





If the available bound is \(H_k/N\le B_k\), multiplication gives \(H_k\le NB_k\). Combining it with parameter data processing yields \(H_k^\theta\le NB_k\), not \(B_k\). The same reasoning explains why the full initial entropy persists in the parameter theorem: the joint evolution was controlled using that initial quantity, and the projection inequality alone gives no replacement of it by a smaller projected initial value.



Earlier results used here: \cite[Section~3.1.1]{LangevinLean2026IPLA}; \cite[Section~3.1.6]{LangevinLean2026IPLA}.
\subsection{\texorpdfstring{Specialize all transfers to the actual Gibbs parameter target}{Specialize all transfers to the actual Gibbs parameter target}}
\label{report-step-IPLA-LSI-Marginal-lean-4}


Set \(\nu=\Pi_N\) in the three generic statements. By the definition of \(\Pi_N^\theta\), this yields extended data processing \(\operatorname{KL}(p_\#\mu\Vert\Pi_N^\theta)\le\operatorname{KL}(\mu\Vert\Pi_N)\), finiteness of parameter KL whenever joint KL is finite, and any real upper-bound transfer. The proofs are direct specializations and retain normalization of both laws and the finite-joint guard. No additional concentration or optimizer estimate appears in this measure projection step.





Both measures in data processing must be projected by the same map. Using the actual Gibbs projection ensures that the reference in the parameter conclusion corresponds exactly to the reference in the joint theorem. The probability and finiteness hypotheses are inherited from that theorem. Consequently the specialization adds no density regularity, independence or optimizer assumption beyond those already used to control the joint law.



Earlier results used here: \cite[Section~3.1.1]{LangevinLean2026IPLA}; \cite[Section~3.1.6]{LangevinLean2026IPLA}; \cite[Section~3.3.2]{LangevinLean2026IPLA}.
\section{\texorpdfstring{Unconditional IPLA entropy from the actual lifted LSI}{Unconditional IPLA entropy from the actual lifted LSI}}
\label{proof-module-IPLA-LSI-UnconditionalTheorem-lean}


This route uses the constructed Gaussian interpolation of the shared entropy theorem. For a positive physical step satisfying the LSI-based restriction and finite initial joint entropy, it derives both finite output entropy and the numerical bound for actual particle laws. Normalization and parameter projection are then algebraic or measure-theoretic consequences of that joint result.




The actual IPLA state laws are ULA laws for the scaled potential \(V_N\), with step \(\gamma/N\), smoothness \(NL\), and dimension \(D=d_\theta+Nd_x\). Applying the shared unconditional ULA theorem proves finite joint KL and its convergence estimate from smoothness, the actual lifted-target LSI, the positive step restriction, and finite initial KL. This also gives normalized, linear-LSI, and parameter-projection consequences.



Earlier results used here: \cite[Section~3.29]{LangevinLean2026ULALSI}.
\subsection{\texorpdfstring{Finite joint entropy and the full IPLA bound}{Finite joint entropy and the full IPLA bound}}
\label{report-step-IPLA-LSI-UnconditionalTheorem-lean-1}


Let \(N>0\), \(D=d_\theta+Nd_x\), and \(V_N(z)=\sum_iU(\theta,\sqrt N z_i)\). Assume differentiable \(U\) with \(L\)-Lipschitz gradient, a genuine LSI with constant \(\beta>0\) for the actual lifted Gibbs target \(\Pi_N\), a positive step \(\gamma\le\beta/(4NL^2)\), and finite initial joint entropy \(H_0\). Then every actual particle law has finite KL and

\[H_k\le e^{-\beta\gamma k/N}H_0+\frac{8\gamma ND L^2}{\beta}.\]

The lifted potential is differentiable and has \(NL\)-Lipschitz gradient. The particle law is exactly its ULA law with step \(h=\gamma/N\). The step equivalence gives \(h\le\beta/(4(NL)^2)\), so the shared unconditional ULA theorem applies. Substituting \(h=\gamma/N\), dimension \(D\), and smoothness \(NL\) transforms its rate into \(e^{-\beta\gamma k/N}\) and its bias into \(8\gamma ND L^2/\beta\). The proof derives positive \(L\) from the positive admissible step before dividing.



The shared theorem constructs all Gaussian interpolation and entropy calculus. Neither compatibility nor regular interpolation data is a premise here. Finiteness is retained as the first component of the conclusion, before real-valued KL is used.




The lifted substitution must be made in both occurrences of smoothness: \(\ell^2=N^2L^2\), while \(h=\gamma/N\). Thus the step restriction becomes \(\gamma\le\beta/(4NL^2)\) and the bias becomes \(8\gamma NDL^2/\beta\). Equality of the Gaussian Euler and particle laws transfers these conclusions to the algorithm. Finite entropy is kept alongside the estimate so every real divergence appearing in the formula is justified.



Earlier results used here: \cite[Section~3.1.3]{LangevinLean2026ULALSI}; \cite[Section~3.29.7]{LangevinLean2026ULALSI}; \cite[Section~3.1.1]{LangevinLean2026IPLA}; \cite[Section~3.1.3]{LangevinLean2026IPLA}; \cite[Section~3.1.4]{LangevinLean2026IPLA}; \cite[Section~3.3.2]{LangevinLean2026IPLA}.
\subsection{\texorpdfstring{Normalize the full joint entropy}{Normalize the full joint entropy}}
\label{report-step-IPLA-LSI-UnconditionalTheorem-lean-2}


Start with the finite full-state entropy theorem. Divide its real inequality by the positive number \(N\), preserving the inequality direction:

\[\frac{H_k}{N}\le e^{-\beta\gamma k/N}\frac{H_0}{N}+\frac{8\gamma D L^2}{\beta}.\]


Since \(N>0\), division preserves the order and gives the same exponential multiplying the intensive initial entropy. The full divergence remains a finite property of the original law; only its numerical normalization has changed. This per-particle quantity is useful for comparing enlarged systems of different size, but it is not itself the entropy of the parameter marginal.



Earlier results used here: \cite[Section~3.1.3]{LangevinLean2026ULALSI}; \cite[Section~3.1.1]{LangevinLean2026IPLA}; \cite[Section~3.3.2]{LangevinLean2026IPLA}.
\subsection{\texorpdfstring{An actual lifted LSI of strength N alpha}{An actual lifted LSI of strength N alpha}}
\label{report-step-IPLA-LSI-UnconditionalTheorem-lean-3}


Suppose the actual lifted target satisfies LSI with constant \(\beta=N\alpha\), where \(\alpha>0\). The positive particle count gives \(\beta>0\), and the step equivalence simplifies \(\gamma\le\beta/(4NL^2)\) to \(\gamma\le\alpha/(4L^2)\). Apply the unconditional normalized theorem at this lifted constant.



The rate simplifies to \(e^{-\alpha\gamma k}\). Since \(D/N=d_x+d_\theta/N\), its normalized bias becomes \((8\gamma L^2/\alpha)(d_x+d_\theta/N)\). Thus

\[\frac{H_k}{N}\le e^{-\alpha\gamma k}\frac{H_0}{N}+\frac{8\gamma L^2}{\alpha}\left(d_x+\frac{d_\theta}{N}\right).\]

The LSI premise is for \(\Pi_N\) itself. The proof performs these exact substitutions and retains finite output KL; it does not infer the lifted LSI by tensorizing a parameter–latent pair model.




The additional LSI premise has to concern the actual coupled particle target. Its substitution cancels the particle factor in the exponential and combines the normalized dimension as \(D/N=d_x+d_\theta/N\). Consequently the intensive floor has a finite latent-dimension term at fixed physical step. The argument performs these cancellations after invoking the constructed joint theorem and does not invoke pairwise product tensorization.



Earlier results used here: \cite[Section~3.1.3]{LangevinLean2026ULALSI}; \cite[Section~3.29.1]{LangevinLean2026ULALSI}; \cite[Section~3.1.1]{LangevinLean2026IPLA}; \cite[Section~3.3.2]{LangevinLean2026IPLA}.
\subsection{\texorpdfstring{Project the actual joint law to the parameter}{Project the actual joint law to the parameter}}
\label{report-step-IPLA-LSI-UnconditionalTheorem-lean-4}


Apply the parameter projection to the actual particle law and to the same lifted Gibbs target. KL data processing gives \(\operatorname{KL}(\mu_k^\theta\Vert\Pi_N^\theta)\le H_k\). The unconditional joint theorem first supplies \(H_k<\infty\), which also proves finiteness of the projected KL before real conversion.



Transitivity with the full joint estimate yields

\[\operatorname{KL}(\mu_k^\theta\Vert\Pi_N^\theta)\le e^{-\beta\gamma k/N}H_0+\frac{8\gamma ND L^2}{\beta}.\]

The initial entropy on the right is the full joint \(H_0\). The bound is unnormalized, so the factor \(N\) in the bias is retained. A per-particle joint estimate must be multiplied by \(N\) before it bounds this unnormalized parameter divergence.




Project both the current law and reference by the same parameter map. The finite joint divergence bounds the projected divergence before the joint numerical estimate is substituted. The initial term remains full joint entropy because it came from the joint evolution. As no per-particle normalization is taken here, the bias retains its full particle factor; dividing it silently would change the quantity being bounded.



Earlier results used here: \cite[Section~3.1.3]{LangevinLean2026ULALSI}; \cite[Section~3.1.1]{LangevinLean2026IPLA}; \cite[Section~3.1.6]{LangevinLean2026IPLA}; \cite[Section~3.3.2]{LangevinLean2026IPLA}.
\section{\texorpdfstring{Conditional entropy bounds for the actual parameter marginal}{Conditional entropy bounds for the actual parameter marginal}}
\label{proof-module-IPLA-LSI-Corollaries-lean}


The parameter conclusions combine two previously established facts: finite entropy and a quantitative bound for the joint Euler law, and monotonicity of divergence under the parameter projection. The complete initial joint entropy and any factor arising from per-particle normalization are retained in every resulting inequality.




Combine the finite joint entropy result with measurable parameter projection. Every conclusion retains the lifted-target LSI and the explicit regular interpolation data. The initialization term remains the full joint initial divergence.

\subsection{\texorpdfstring{The actual parameter divergence is finite}{The actual parameter divergence is finite}}
\label{report-step-IPLA-LSI-Corollaries-lean-1}


The particle-law probability theorem proves \(\rho_k\) is a probability, and the LSI hypothesis proves \(\Pi_N\) is a probability. The supplied interpolation data prove finite actual joint KL for \(\gamma\ge0\). Apply the Gibbs-specific finiteness transfer to obtain \(\operatorname{KL}(p_{\theta\#}\rho_k\Vert\Pi_N^\theta)<\infty\). This is an extended-valued assertion, proved before any real-valued inequality.





The finite joint entropy inequality already places the current law in the finite-divergence domain relative to the target. Data processing maps this domain into the corresponding parameter finite-divergence domain. This proves the existence of a finite parameter entropy before assigning a numerical upper bound to it. Probability normalization of both laws is inherited from the actual Markov transition and the stated target LSI.



Earlier results used here: \cite[Section~3.1.3]{LangevinLean2026ULALSI}; \cite[Section~3.34.1]{LangevinLean2026ULALSI}; \cite[Section~3.1.1]{LangevinLean2026IPLA}; \cite[Section~3.1.6]{LangevinLean2026IPLA}; \cite[Section~3.3.2]{LangevinLean2026IPLA}.
\subsection{\texorpdfstring{Transfer the full general joint bound}{Transfer the full general joint bound}}
\label{report-step-IPLA-LSI-Corollaries-lean-2}


Under \(N>0\), \(\beta>0\), \(\beta\le NL\), and \(0<\gamma\le\beta/(4NL^2)\), the joint theorem gives \(H_k\le e^{-\beta\gamma k/N}H_0+8\gamma ND L^2/\beta\). The finite joint guard and the Gibbs projection bound yield

\[\operatorname{KL}(p_{\theta\#}\rho_k\Vert\Pi_N^\theta)\le e^{-\beta\gamma k/N}H_0+\frac{8\gamma ND L^2}{\beta}.\]

Both probability properties are established within the proof. Here \(H_0=\operatorname{KL}(\mu_0\Vert\Pi_N)\), rather than the smaller projected initial divergence; data processing does not permit replacing the initial term in this argument.





Apply data processing at the same iteration \(k\), then substitute the full joint bound. Its initial term measures information in both the parameter and particles. The exponential and bias are unchanged because projection only decreases the left side. This argument controls divergence between parameter probability laws; it does not convert that divergence into distance from a deterministic maximizing parameter.



Earlier results used here: \cite[Section~3.1.3]{LangevinLean2026ULALSI}; \cite[Section~3.34.1]{LangevinLean2026ULALSI}; \cite[Section~3.1.1]{LangevinLean2026IPLA}; \cite[Section~3.1.6]{LangevinLean2026IPLA}; \cite[Section~3.3.2]{LangevinLean2026IPLA}.
\subsection{\texorpdfstring{Retain N when using the normalized linear-\allowbreak{}LSI bound}{Retain N when using the normalized linear-LSI bound}}
\label{report-step-IPLA-LSI-Corollaries-lean-3}


If the actual lifted LSI constant is \(N\alpha\), the normalized joint theorem gives

\[\frac{H_k}{N}\le e^{-\alpha\gamma k}\frac{H_0}{N}+\frac{8\gamma L^2}{\alpha}\left(d_x+\frac{d_\theta}{N}\right).\]

The per-particle transfer lemma multiplies this complete upper bound by positive \(N\), so the unscaled parameter divergence is at most

\[N\left[e^{-\alpha\gamma k}\frac{H_0}{N}+\frac{8\gamma L^2}{\alpha}\left(d_x+\frac{d_\theta}{N}\right)\right].\]

The proof obtains finite joint KL separately and then applies this guarded transfer. This entropy estimate compares two parameter probability laws; no concentration at a marginal-likelihood maximizer is asserted.




The proof first installs probability instances for the actual particle law and Gibbs target, using kernel mass preservation and the LSI hypothesis's normalization field. It next supplies the proved finite joint KL to the real data-processing theorem. In the general case it transfers the full joint bound directly. In the \(N\alpha\) case it transfers the normalized bound by multiplying its entire right side by \(N\): \[\operatorname{KL}(p_\#\rho_k\Vert p_\#\Pi_N)\le N\left[e^{-\alpha\gamma k}\frac{H_0}{N}+\frac{8\gamma L^2}{\alpha}(d_x+d_\theta/N)\right].\] The lifted LSI and regular interpolation hypotheses remain exactly those of the joint theorem.





Under the linear-LSI specialization, start from the intensive estimate and undo its normalization before projecting. Multiplying the complete right side by \(N\) restores the initial entropy \(H_0\) and produces bias \(8\gamma L^2D/\alpha\). This is consistent with substituting \(\beta=N\alpha\) directly into the full theorem. Agreement of the two routes checks the particle factors in the corollary.



Earlier results used here: \cite[Section~3.1.3]{LangevinLean2026ULALSI}; \cite[Section~3.34.1]{LangevinLean2026ULALSI}; \cite[Section~3.1.1]{LangevinLean2026IPLA}; \cite[Section~3.1.6]{LangevinLean2026IPLA}; \cite[Section~3.3.2]{LangevinLean2026IPLA}.

\begin{thebibliography}{ULA-LSI}
\bibitem[ULA-SC]{LangevinLean2026ULASC}
LangevinLean Project. \emph{Unadjusted Langevin under Strong Convexity}. LL-01, 2026. Draft report, content revision 8d6335dce333. \url{https://odakyildiz.github.io/LangevinLean/docs/pdf/UnadjustedLangevin/FullProof.pdf}.
\bibitem[ULA-LSI]{LangevinLean2026ULALSI}
LangevinLean Project. \emph{Unadjusted Langevin under LSI}. LL-02, 2026. Draft report, content revision e64aa62fb836. \url{https://odakyildiz.github.io/LangevinLean/docs/pdf/UnadjustedLangevin/LSI/FullProof.pdf}.
\bibitem[IPLA-SC]{LangevinLean2026IPLA}
LangevinLean Project. \emph{Interacting Particle Langevin under Strong Convexity}. LL-03, 2026. Draft report, content revision c72e54460954. \url{https://odakyildiz.github.io/LangevinLean/docs/pdf/IPLA/FullProof.pdf}.
\end{thebibliography}

\end{document}
